\subsection{Topological sphere concordances and finite-rank stabilization}
We express ordinary concordances of spheres by the failure of
finite-rank smoothing spaces to stabilize. The proof first uses
parametric germ extension to obtain a cone fibre, then identifies
its map with product by a line.


Let $\mathrm{TOP}(k)=\operatorname{Homeo}(\mathbb R^k;0)$ carry its
actual compact-open topology, and put
\begin{equation}\label{eq:finite-rank-top-smoothing-space}
 Y_k=\operatorname{hofib}_{*}
       \bigl(B\mathrm O(k)\longrightarrow B\mathrm{TOP}(k)\bigr).
\end{equation}
The classifying spaces may be formed from singular simplicial groups.
The orthogonal group has its usual topology, and its inclusion in
$\mathrm{TOP}(k)$ is the literal linear inclusion.  Product with a
line gives a based stabilization $Y_k\to Y_{k+1}$.

\begin{theorem}\label{thm:top-sphere-finite-rank-stabilization}
For every standard sphere $S^n$, $n\geq2$, the ordinary concordance
space with bottom fixed pointwise and top free satisfies
\begin{equation}\label{eq:top-sphere-concordance-fibre}
 P_{\mathrm{TOP}}(S^n)\simeq
       \operatorname{hofib}_{*}(Y_n\longrightarrow Y_{n+1}).
\end{equation}
These are the finite-rank spaces in
\eqref{eq:finite-rank-top-smoothing-space}.
\end{theorem}

\begin{proof}
We work with singular simplicial groups for actual homeomorphism
groups.  Write $\mathcal G_k$ for the simplicial group of germs at
the zero section: a simplex is represented on
$\Delta^q\times U$, for one open neighbourhood $U$ of $0$, by a
projection-preserving homeomorphism onto an open image that fixes
$\Delta^q\times\{0\}$.  Two representatives agree as germs when
they agree on a smaller common neighbourhood for the whole simplex.
This is the finite-parameter convention of
\cite[Section~3.4, equations~(3.5)--(3.7), pp.~14--15]
{SalvatoreTurchin2026}.

\medskip\noindent\textit{Global representatives and point stabilizers.}
Set $T_k=\operatorname{Sing}\mathrm{TOP}(k)$.  Restriction induces
a degreewise surjection
\[
                     T_k\longrightarrow\mathcal G_k.
\]
Indeed, a germ family has a common source ball and a common image
chart by compactness of its parameter simplex.  After fixed changes
of scale, Kuiper--Lashof's parametric extension lemma gives a
sphere homeomorphism family agreeing with it on a smaller ball and
fixing both distinguished points
\cite[Lemma~3, p.~11; TOP proof, pp.~12--13]{KuiperLashof1966I}.
Deleting the point at infinity gives the required global family.

A degreewise surjective map of simplicial groups is a Kan fibration.
Its kernel consists of families equal to the identity on a common
ball about $0$.  Sphere inversion identifies this kernel with
families of compactly supported homeomorphisms of $\mathbb R^k$
that fix $0$.  Every finite family has a common bounded support.
The formula
\[
 a_s(x)=s\,a(x/s)\quad(s>0),\qquad a_0(x)=x
\]
contracts such a family, together with its inverses, in the
compact-open topology.  For example, if the common support lies
in $D_R$, then the support at time $s$ lies in $D_{sR}$ and
$\sup\|a_s-\operatorname{id}\|\leq
s\sup\|a-\operatorname{id}\|$.  Identity parameter values stay
fixed.  Thus the kernel is weakly contractible and
$T_k\to\mathcal G_k$ is a weak equivalence on all components.

Fix $p\in S^n$ and an $\mathrm O(n)$-equivariant chart at $p$.
For $H_n=\operatorname{Sing}\operatorname{Homeo}(S^n;p)$ the
germ map $H_n\to\mathcal G_n$ is also degreewise surjective by
the same lemma.  Its kernel is the space of families equal to the
identity on a common neighbourhood of $p$.  On the complement of
$p$ these are compactly supported families on $\mathbb R^n$.
The same contraction works, without any additional fixed-point
condition.  Hence
\begin{equation}\label{eq:sphere-stabilizer-germ-equivalence}
 H_n\xrightarrow{\simeq}\mathcal G_n
                 \xleftarrow{\simeq}T_n.
\end{equation}
Both maps are the identity on the continuous orthogonal singular
subgroup specified by the chosen chart.

\medskip\noindent\textit{The disk germ fibre.}
Let $m=n+1$ and
$A_m=\operatorname{Sing}\operatorname{Homeo}(D^m;0)$.
The germ map $A_m\to\mathcal G_m$ is a Kan fibration.  We recall
the relative argument, which has no restriction on $m$.
The Edwards--Kirby local ambient-extension operator assigns to
every sufficiently small perturbation $v$ of the inclusion of a
compact ball $U\subset\operatorname{int}D^m$ a disk homeomorphism
$\sigma(v)$ equal to $v$ on a smaller ball $C$, supported away
from the outer boundary, with $\sigma(\operatorname{id})=
\operatorname{id}$
\cite[Theorem~5.1, p.~70; proof of Corollary~1.2, p.~79]
{EdwardsKirby1971}.  It depends continuously on $v$.

Suppose a germ simplex $f_t$ has a disk lift $a_t$ on a horn
$L\subset\Delta^q$.  Choose a deformation $d_s$ of the simplex
from a retraction onto $L$ to the identity, fixed on $L$.  On
common smaller neighbourhoods put
\[
 a^0_t=a_{d_0(t)},\qquad
 k_{t,s}=f_{d_s(t)}(a^0_t)^{-1}.
\]
The equalities $k_{t,0}=1$ and $k_{t,s}=1$ for $t\in L$ hold
on one common ball, after shrinking once.  Choose $U,C$ and a
still smaller ball $V$ so that all increments
$k_{t,s}k_{t,s'}^{-1}$ are defined on $U$ and
$k_{t,s}(V)\subset\operatorname{int}C$ for every $(t,s)$.
Joint continuity of maps and inverses and compactness give these
choices.  They also give a single finite partition of $I$ on
whose subintervals every increment belongs to the domain of
$\sigma$.  Successively multiplying their extensions produces
$b_{t,s}$ with germ $k_{t,s}$, equal to $1$ on the horn.  The
induction is exact on $V$, since the previous image of $V$ lies
in $C$.  Then $b_{t,1}a^0_t$ is a lift equal to $a_t$ on $L$.
Every extension fixes $0$ because the corresponding increment
does so near $0$.  This proves the fibration assertion with the
prescribed faces fixed exactly.

Its identity fibre is the filtered union of
\[
 K_\varepsilon=
 \operatorname{Sing}\operatorname{Homeo}
                  (D^m\;\mathrm{rel}\;D^m_\varepsilon),
 \qquad 0<\varepsilon<1.
\]
Restriction to the annulus identifies each $K_\varepsilon$ with
$P_{\mathrm{TOP}}(S^n)$: the inner boundary is fixed and the
outer boundary is free.  If $\varepsilon<\varepsilon_0$, radial
conjugations that move $D_{\varepsilon_0}$ to $D_\varepsilon$
give a homotopy in $K_\varepsilon$ from the inclusion
$K_{\varepsilon_0}\to K_\varepsilon$ to an isomorphism.
Each finite simplicial sphere or homotopy in the union has one
common fixed ball.  Therefore the fibre has the weak homotopy
type of $P_{\mathrm{TOP}}(S^n)$.

Write $G_n=\operatorname{Sing}\operatorname{Homeo}(S^n)$.
Boundary restriction $A_m\to G_n$ is a weak equivalence: radial
coning gives a section, and the boundary-fixed kernel, still
fixing $0$, contracts by the relative Alexander formula.
Extend a disk map outside the disk by coning its boundary map.
This gives a homomorphism $A_m\to T_m$ compatible with germs.
On the cone section it is the radial homomorphism
\[
 c:G_n\longrightarrow T_{n+1},\qquad
                  c(h)(r\omega)=r h(\omega).
\]
The preceding weak equivalences and the disk germ fibration give
\begin{equation}\label{eq:top-sphere-cone-fibre}
             P_{\mathrm{TOP}}(S^n)\simeq\operatorname{hofib}_1(c).
\end{equation}
The global germ equivalence here includes all components.

\medskip\noindent\textit{Removing orthogonal directions.}
The cone map is the linear inclusion on
$\operatorname{Sing}\mathrm O(n+1)$.  Taking right cosets by
this same simplicial subgroup preserves its homotopy fibre:
\[
 \operatorname{hofib}(c)\simeq
 \operatorname{hofib}\left(
   G_n/\operatorname{Sing}\mathrm O(n+1)
       \longrightarrow T_{n+1}/\operatorname{Sing}\mathrm O(n+1)
                         \right).
\]
This follows from the two simplicial subgroup fibrations, whose
map on fibres is the identity.  Such a quotient models the
homotopy fibre of the corresponding map of classifying spaces.
We use these simplicial subgroup quotients throughout.

The point stabilizer gives
\[
 H_n/\operatorname{Sing}\mathrm O(n)
       \xrightarrow{\simeq}
 G_n/\operatorname{Sing}\mathrm O(n+1).
\]
To see this directly, for a family $g_t$ on a simplex choose
rotations $R_t$ taking $p$ to $g_t^{-1}(p)$.  Such a continuous
choice exists because the pullback of
$\mathrm O(n+1)\to S^n$ to a simplex is a trivial principal
$\mathrm O(n)$-bundle.  Then $g_tR_t$ fixes $p$, and two such
representatives differ precisely by a family in $\mathrm O(n)$.
Together with \eqref{eq:sphere-stabilizer-germ-equivalence},
this identifies the source quotient with $Y_n$ and the target
quotient with $Y_{n+1}$.

\medskip\noindent\textit{Identifying the arrow.}
It remains to identify this map, not just its source and target.
Take $p=(0,\ldots,0,1)$.  If $h$ fixes $p$, its cone fixes the
whole ray through $p$.  Consequently
\[
 F_s(h)=\tau_{-sp}\,c(h)\,\tau_{sp},\qquad 0\leq s\leq1,
\]
is a homotopy of homomorphisms from $H_n$ to $T_{n+1}$, all
fixing $0$.  It is constant on each
$\operatorname{diag}(A,1)$ with $A\in\mathrm O(n)$.
Near the origin use
\[
 \theta(u)=(u,\sqrt{1-\|u\|^2}),\qquad
 \psi(u,t)=(1+t)\theta(u)-p.
\]
The cone preserves radius and applies $h$ to direction.  Thus
the coordinates centred at $p$ give the exact germ identity
\begin{equation}\label{eq:cone-germ-line-stabilization}
 \psi^{-1}F_1(h)\psi(u,t)
               =\bigl(\theta^{-1}h\theta(u),t\bigr).
\end{equation}
The right side is product with a line on the point germ of $h$.

The coordinate change does not alter the homotopy class of this
map.  For $r>0$, let $\psi_r(z)=r^{-1}\psi(rz)$ and set
$\psi_0=\operatorname{id}$.  These form a continuous family
of local diffeomorphisms with continuous inverses on common
neighbourhoods.  For example, writing $y=(v,w)$ and
$N_r=\|(rv,1+rw)\|$, the inverse has the nonsingular expression
\[
 \psi_r^{-1}(v,w)=
 \left(\frac{v}{N_r},\,
       \frac{2w+r(\|v\|^2+w^2)}{N_r+1}\right),
 \qquad N_0=1.
\]
This verifies the extension at $r=0$ directly.  The family
commutes with $\operatorname{diag}(A,1)$ for every
$A\in\mathrm O(n)$.  Conjugation by it therefore provides
the required homotopy of quotient maps, relative to the source
orthogonal subgroup.  A finite parameter family of $h$ uses
one smaller chart; identity parameter faces remain fixed.
Equation~\eqref{eq:cone-germ-line-stabilization} now identifies
the arrow in \eqref{eq:top-sphere-cone-fibre} with the actual
stabilization $Y_n\to Y_{n+1}$, proving the theorem.
\end{proof}

\begin{corollary}\label{cor:four-sphere-low-concordance-groups}
For the standard four-sphere,
\[
 \pi_iP_{\mathrm{TOP}}(S^4)=\pi_iP_{\mathrm{PL}}(S^4)=0
                       \qquad(0\leq i\leq3).
\]
For $i=0$, both component sets are singletons.
\end{corollary}

\begin{proof}
The stabilization $Y_4\to\mathrm{TOP}/\mathrm O$ is
$5$-connected by
\cite[Theorem~8.7A, p.~128]{FreedmanQuinn1990}, with the
corrected account in
\cite[Section~1.1, Theorem~1.3 and the following paragraph,
p.~2]{GabaiGayHartmanKrushkalPowell2026}.
The latter uses precisely the global origin-fixing group and
classifying-space fibre convention in
\eqref{eq:finite-rank-top-smoothing-space}.

For rank five, use the direct finite-rank statement
\cite[Annex~C, Theorem~5.3, p.~317]{KirbySiebenmann1977}:
\[
 \pi_i(\mathrm{TOP}_m/\mathrm{CAT}_m)
       \xrightarrow{\cong}\pi_i(\mathrm{TOP}/\mathrm{CAT})
       \quad(i<m,\ m\geq5),
       \qquad \mathrm{CAT}=\mathrm{PL}\text{ or }\mathrm{DIFF}.
\]
The source defines these as classifying-space fibres; its DIFF
version uses the orthogonal classifying space.  The orthogonally compatible germ
comparison above therefore identifies its rank-five map with
$Y_5\to\mathrm{TOP}/\mathrm O$.  It is an isomorphism on
$\pi_i$ for $i\leq4$.  Naturality of stabilization shows that
$Y_4\to Y_5$ is also an isomorphism in those degrees.  The
homotopy-fibre sequence of the theorem gives the TOP vanishing
through degree three, including the component assertion.

For PL use \eqref{eq:sphere-concordance-fiber} with $n=4$.
Here $X_j=\mathrm{PL}_j/\mathrm O_j$ denotes the classical
unstable smoothing homotopy type.  Its classifying-space
definition does not assert a literal inclusion of continuous
orthogonal singular chains into joint-PL germ chains; a precise
construction via piecewise differentiable germs is recalled in
\cite[Section~3.4, Remark~3.7 and footnote~16, p.~16]
{SalvatoreTurchin2026}.  The classical stabilization calculation
\cite[Section~5, formula~(3), p.~37]{BurgheleaLashof1974II}
and the stable vanishing
\cite[Annex~C, p.~318]{KirbySiebenmann1977} give
\[
              \pi_iX_4=0\ (i\leq3),\qquad
              \pi_iX_5=0\ (i\leq4).
\]
The exact sequence of $\operatorname{hofib}(X_4\to X_5)$
proves the PL assertion.  Its source is connected and its
target is simply connected, so the same conclusion includes
components.
\end{proof}

The argument does not determine $\pi_4P_{\mathrm{TOP}}(S^4)$:
the displayed stabilization information does not control
$\pi_5Y_4\to\pi_5Y_5$.  It also leaves the corresponding
ordinary concordance spaces for general four-manifold bases
unevaluated.
